\documentclass[11pt]{amsart}
\usepackage{amsmath,amssymb,amsthm,booktabs}
\usepackage[margin=1.15in]{geometry}
\usepackage{hyperref}
\newtheorem{theorem}{Theorem}
\newtheorem{proposition}[theorem]{Proposition}
\newtheorem{lemma}[theorem]{Lemma}
\newtheorem{corollary}[theorem]{Corollary}
\newtheorem{conjecture}[theorem]{Conjecture}
\theoremstyle{remark}
\newtheorem{remark}[theorem]{Remark}
\newtheorem{example}[theorem]{Example}
\newcommand{\PP}{\mathbb{P}}
\newcommand{\EE}{\mathbb{E}}
\newcommand{\cF}{\mathcal{F}}
\newcommand{\one}{\mathbf{1}}

\title{Dead branches in the trace-reinforced ant process}
\author{Yitzchak Shmalo}
\date{November 2026}

\begin{document}
\begin{abstract}
In the trace-reinforced ant process of Kious, Mailler and Schapira, ants walk one after another from a nest $N$ to a
food source $F$ on a finite graph, and each ant adds one unit of weight to every edge it has crossed. It was conjectured
that the normalised weights converge to a deterministic limit, and that this limit is positive on every edge as soon as
$N$ and $F$ are at distance at least two. We show that the second statement fails. If a part of the graph hangs by a
single edge from a neighbour of $F$, on the side away from $N$, then the normalised weights of all its edges tend to
zero almost surely; this happens in $2$-connected graphs with $N$ and $F$ at distance two. We also show that the
mean-field vector field has no equilibrium with all coordinates positive when some neighbour of $F$ is hidden from $N$
by the other neighbours of $F$, so that on such graphs, some of which are $3$-connected, the normalised weights cannot
converge to a positive limit. The equilibria are described through the voltage of the unit current from $N$ to $F$.
A census of the $22{,}828$ cases on connected graphs with at most seven vertices supports a corrected form of the
conjecture, in which the set of edges with limit zero is given by an explicit combinatorial rule. We also give the
equilibrium on a graph made of any number of parallel paths.
\end{abstract}
\maketitle

\section{Introduction}

Let $G=(V,E)$ be a finite connected simple graph with two distinct marked vertices, the nest $N$ and the food $F$.
The \emph{trace-reinforced ant process} on $(G,N,F)$ is the sequence of weight vectors $W(n)=(W_e(n))_{e\in E}$
defined as follows. Set $W_e(0)=1$ for all $e$. Given $W(n)$, the $(n+1)$-st ant performs a random walk started at
$N$ which jumps from a vertex $x$ to a neighbour $z$ with probability $W_{xz}(n)/W_x(n)$, where
$W_x(n)=\sum_{z\sim x}W_{xz}(n)$, and stops when it first hits $F$. Then
\[
W_e(n+1)=W_e(n)+\one\{\text{the $(n+1)$-st ant crossed $e$}\}.
\]
The process was introduced by Kious, Mailler and Schapira \cite{KMS}, as a variant of the models of \cite{KMS0}, in
which an ant reinforces only a part of its trace (its loop-erasure, or a shortest path inside it). These models
are motivated by the trails with which ants mark the way to a food source \cite{Den,DS}, and belong to the family of
reinforced random processes surveyed in \cite{Pem07}, which goes back to the edge-reinforced walk \cite{Dav}; a model with two nests is treated in \cite{MV}. Kious, Mailler
and Schapira proved that $W(n)/n$ converges almost surely to an explicit
deterministic limit on several families of graphs (tree-like graphs in which $N$ and $F$ are adjacent, the cone, two
parallel paths, the lozenge), and stated the following.

\begin{conjecture}[\cite{KMS}, Conjecture 1.1]\label{conj:kms}
There exists a deterministic family $(\chi_e)_{e\in E}$ such that $W_e(n)/n\to\chi_e$ almost surely for all $e$.
Furthermore, if in addition the distance between $N$ and $F$ is at least $2$, then $\chi_e>0$ for all $e\in E$.
\end{conjecture}

An edge with an endpoint outside the connected component of $N$ in $G-F$ is never crossed, so the second sentence is
to be read for graphs in which $G-F$ is connected; we call such triples $(G,N,F)$ \emph{reduced}. The second sentence
fails for reduced graphs.

\begin{theorem}[Dead branches]\label{thm:dead}
Let $u$ be a neighbour of $F$, and let $H\subset V\setminus\{N,F,u\}$ be a set of vertices such that exactly one edge
$e=\{u,y\}$ joins $H$ to $V\setminus(H\cup\{F\})$. Then, almost surely,
\[
\frac{W_f(n)}{n}\longrightarrow 0\qquad\text{for every edge $f$ with an endpoint in $H$.}
\]
\end{theorem}

In other words, $H$ is the far side of a bridge of $G-F$ attached at a neighbour of $F$; the vertices of $H$ may have
any number of edges to $F$. The vertex $u$ may be $N$ itself. When $G-F$ is a tree and $u=N$ is adjacent to $F$ the
theorem covers every edge other than $\{N,F\}$, which is the second half of Theorem 1.3 of \cite{KMS}. When $N$ is
adjacent to $F$, Proposition 2.8 of \cite{KMS} shows that $W_{NF}(n)/n\to1$ and that the normalised weight of every
other edge at $F$ tends to zero; Theorem~\ref{thm:dead} needs no such assumption.

\begin{example}\label{ex:paw}
Take a triangle $x,y,F$ and attach $N$ to $x$ by an edge. The distance from $N$ to $F$ is $2$, and
Theorem~\ref{thm:dead} with $u=x$, $H=\{y\}$ gives $W_{xy}(n)/n\to0$ and $W_{yF}(n)/n\to0$. This case also follows
from \cite[Theorem 1.3]{KMS} applied with nest $x$: every ant goes through $x$, and its excursions from $x$ to the
pendant vertex $N$ do not change which of the other edges it crosses. The same reduction applies whenever a
cut vertex separates $N$ from $F$, so the examples of interest are those without one.
\end{example}

\begin{example}\label{ex:c4}
Take the $4$-cycle $N,a,F,b$ and add a vertex $y$ joined to $a$ and to $F$. The graph is $2$-connected, $N$ and $F$
are at distance $2$, and every edge lies on a simple path from $N$ to $F$. Theorem~\ref{thm:dead} with $u=a$,
$H=\{y\}$ gives $W_{ay}(n)/n\to0$ and $W_{yF}(n)/n\to0$: the detour $a,y,F$ dies next to the edge $\{a,F\}$.
\end{example}

The second result concerns the mean-field vector field. For a weight vector $w\in[0,\infty)^E$ whose support connects
$N$ to $F$, let $p_e(w)$ be the probability that the $w$-weighted random walk started at $N$ crosses $e$ before
hitting $F$. Then $\EE[W_e(n+1)-W_e(n)\mid\cF_n]=p_e(W(n))$, the map $p$ is homogeneous of degree $0$, and the process
is a stochastic approximation of the flow $\dot w=p(w)-w$ in the sense of \cite{Ben,Pem90}. An \emph{equilibrium} is a solution of $p(w)=w$.

\begin{theorem}[Shadowed neighbours of $F$]\label{thm:shadow}
Let $(G,N,F)$ be reduced with $N$ not adjacent to $F$, let $N(F)$ be the set of neighbours of $F$, and let $I_0$ be
the connected component of $N$ in $G-F-N(F)$. If some $x\in N(F)$ has no neighbour in $I_0$, then $p$ has no
equilibrium with all coordinates positive. Consequently, on such a graph the event that $W(n)/n$ converges to a limit
with all coordinates positive has probability zero.
\end{theorem}

So on these graphs the two sentences of Conjecture~\ref{conj:kms} cannot both hold: either $W(n)/n$ does not converge,
or its limit has a zero coordinate. The hypothesis does not involve bridges or cut vertices. For instance, take
the complete bipartite graph with parts $\{a,b,c\}$ and $\{N,F,x\}$ and add the edge $\{x,F\}$. This graph is
$3$-connected, $N$ and $F$ are at distance $2$, $I_0=\{N\}$, and $x$ is a neighbour of $F$ with no neighbour in $I_0$.
A smaller example is the $4$-cycle $N,a,F,b$ with a fifth vertex $c$ joined to $a$, $b$ and $F$. In the mean-field
flow on this graph the edge $\{c,F\}$ dies while $\{a,c\}$ and $\{b,c\}$ keep the positive weight $0.3209$
(Section~\ref{sec:census}).

Section~\ref{sec:proof} contains the proof of Theorem~\ref{thm:dead}. Section~\ref{sec:eq} describes the equilibria
through electrical voltages and proves Theorem~\ref{thm:shadow}. Section~\ref{sec:census} states a corrected form of
the second sentence of Conjecture~\ref{conj:kms} and reports a census of small graphs. Section~\ref{sec:theta}
computes an equilibrium on theta graphs with any number of paths.

\section{Proof of Theorem~\ref{thm:dead}}\label{sec:proof}

Write $a=\{u,F\}$, $X_n=W_e(n)$ and $A_n=W_a(n)$, and let $\cF_n$ be the $\sigma$-field generated by the first $n$
ants.

\medskip\noindent\textbf{Step 0.} \emph{For every edge $f$ with an endpoint in $H$ and every $n$, $W_f(n)\le X_n$.}
An ant starts at $N\notin H$ and stops at $F$, and $e$ is the only edge from $H$ to $V\setminus(H\cup\{F\})$. Hence an
ant that crosses $f$ has crossed $e$ earlier. Both weights start at $1$ and each increases by one when its edge is
crossed.

\medskip\noindent\textbf{Step 1.} \emph{One ant.} Fix $n$ and write $w=W(n)$, $X=X_n$, $A=A_n$, and $R$ for the
total weight of the edges at $u$ other than $a$ and $e$. Let $h$ be the probability that the walk from $N$ hits $u$
before $F$ ($h=1$ if $u=N$). The walk cannot enter $H$ before it visits $u$. After the first visit to $u$, decompose
the walk into its successive departures from $u$. A departure is of one of four kinds:
\begin{itemize}
\item[(i)] the step $u\to F$, with probability $A/w_u$; the edge $a$ is crossed and the walk ends;
\item[(ii)] the step $u\to y$, with probability $X/w_u$; the edge $e$ is crossed; afterwards the walk returns to $u$
with some probability $\rho$ and is absorbed at $F$ from inside $H$ with probability $1-\rho$;
\item[(iii)] a step along another edge, with probability $R/w_u$, after which the walk returns to $u$ before hitting
$F$; this has conditional probability $\rho'$, say;
\item[(iv)] a step along another edge after which the walk hits $F$ before returning to $u$.
\end{itemize}
An excursion of kind (iii) or (iv) does not meet $H$. Successive departures are independent. The edge $e$ is crossed
if and only if the first departure that is not of kind (iii) is of kind (ii), and the edge $a$ is crossed if and only
if the walk ends by a departure of kind (i). With $S=A+R(1-\rho')$ this gives
\begin{equation}\label{eq:pe}
p_e(w)=\frac{hX}{X+S},\qquad p_a(w)=\frac{hA}{S+X(1-\rho)} .
\end{equation}
Therefore
\begin{equation}\label{eq:diff}
\frac{p_a(w)}{A}-\frac{p_e(w)}{X}=\frac{hX\rho}{(S+X(1-\rho))(S+X)}\ \ge\ \frac{\rho\,p_e(w)}{S+X}.
\end{equation}
Now $S+X\le w_u\le\deg(u)\,(n+1)$, since every weight is at most $n+1$ at time $n$. Moreover $\rho$ is at least the
probability that the first step from $y$ goes back to $u$, which is $X/w_y\ge1/\deg(y)$ by Step 0, all edges at $y$
having an endpoint in $H$. With $c=1/(\deg(u)\deg(y))$ we obtain
\begin{equation}\label{eq:drift}
\frac{p_a(W(n))}{A_n}-\frac{p_e(W(n))}{X_n}\ \ge\ \frac{c\,p_e(W(n))}{n+1}.
\end{equation}

\medskip\noindent\textbf{Step 2.} \emph{The logarithm of the ratio.} Let $L_n=\log(X_n/A_n)$, and let $D_n$ and
$D'_n$ be the indicators that the $(n+1)$-st ant crosses $e$ and $a$ respectively, so that
\[
L_{n+1}-L_n=D_n\log\Bigl(1+\frac1{X_n}\Bigr)-D'_n\log\Bigl(1+\frac1{A_n}\Bigr).
\]
Using $t-t^2/2\le\log(1+t)\le t$ and \eqref{eq:drift},
\[
\EE[L_{n+1}-L_n\mid\cF_n]\ \le\ \frac{p_e(W(n))}{X_n}-\frac{p_a(W(n))}{A_n}+\frac{p_a(W(n))}{2A_n^2}
\ \le\ -\frac{c\,p_e(W(n))}{n+1}+\frac{p_a(W(n))}{2A_n^2}.
\]
Since $X$ increases by one each time $D_n=1$, we have $\sum_nD_n/X_n^2\le\sum_{k\ge1}k^{-2}$, and likewise for $A$.
Hence $\sum_n\EE[(L_{n+1}-L_n)^2]\le2\sum_kk^{-2}<\infty$, so the martingale
$M_n=\sum_{j<n}\bigl(L_{j+1}-L_j-\EE[L_{j+1}-L_j\mid\cF_j]\bigr)$ is bounded in $L^2$ and converges almost surely.
Also $\EE\sum_np_a(W(n))/A_n^2=\EE\sum_nD'_n/A_n^2<\infty$, so $Q=\sum_np_a(W(n))/(2A_n^2)$ is almost surely finite.
Altogether, for all $n$,
\begin{equation}\label{eq:L}
L_n\ \le\ M_n+Q-c\sum_{j<n}\frac{p_e(W(j))}{j+1}.
\end{equation}

\medskip\noindent\textbf{Step 3.} \emph{Conclusion.} The martingale $\sum_{j<n}(D_j-p_e(W(j)))/(j+1)$ has
square-summable increments and converges almost surely. Hence, almost surely,
\begin{equation}\label{eq:equiv}
\sum_j\frac{p_e(W(j))}{j+1}=\infty\quad\Longleftrightarrow\quad\sum_j\frac{D_j}{j+1}=\infty .
\end{equation}
Summation by parts gives $\sum_{j<n}D_j/(j+1)=(X_n-1)/n+\sum_{j=1}^{n-1}(X_j-1)/(j(j+1))$. Suppose that
$X_m\ge\delta m$ for infinitely many $m$, for some $\delta>0$. For such $m$ with $\delta m\ge2$, monotonicity gives
$X_j-1\ge\delta m/2$ for $m\le j\le2m$, and $\sum_{j=m}^{2m}1/(j(j+1))\ge1/(2m)$, so the block $[m,2m]$ contributes at
least $\delta/4$. Infinitely many disjoint blocks give $\sum_jD_j/(j+1)=\infty$. By \eqref{eq:equiv} and \eqref{eq:L},
on the event $\{\limsup X_n/n>0\}$ we have $L_n\to-\infty$ almost surely, that is $X_n/A_n\to0$. But $A_n\le n+1$, so
$X_n/n\to0$ on that event, which is a contradiction unless the event is null. Thus $X_n/n\to0$ almost surely, and
Step~0 extends this to every edge with an endpoint in $H$. \qed

\begin{remark}
The bridge is essential. If the part hanging from $u$ is attached by two edges, say a triangle $u,y_1,y_2$, an ant
can enter by one edge and leave by the other, and the same computation gives $p_{uy_1}(w)/w_{uy_1}>p_a(w)/w_a$ for
small weights on the triangle. In the mean-field flow such a triangle keeps positive weight.
\end{remark}

\section{Equilibria and voltages}\label{sec:eq}

Let $w\in[0,\infty)^E$ have a support that connects $N$ to $F$, and let $g(x)$ be the expected number of visits to
$x$ by the $w$-walk from $N$ before it hits $F$. For a vertex $x\ne F$ with $w_x>0$ put $v(x)=g(x)/w_x$, and put
$v(F)=0$. Then $v$ is the voltage of the unit current from $N$ to $F$ in the network $w$:
\begin{equation}\label{eq:kirchhoff}
\sum_{z\sim x}w_{xz}\,(v(x)-v(z))=\one\{x=N\}\qquad(x\ne F,\ w_x>0).
\end{equation}

\begin{lemma}\label{lem:cross}
The expected number of crossings of $e=\{x,z\}$ by the walk is $w_e\,(v(x)+v(z))$. In particular
$p_e(w)\le w_e(v(x)+v(z))$, and $p_{xF}(w)=w_{xF}\,v(x)$ for every neighbour $x$ of $F$.
\end{lemma}
\begin{proof}
The expected number of steps from $x$ to $z$ is $g(x)\,w_e/w_x$. An edge at $F$ is crossed at most once.
\end{proof}

\begin{proposition}\label{prop:volt}
Let $w$ be an equilibrium and let $V_1$ be the set of neighbours $x$ of $F$ with $w_{xF}>0$. Then
\begin{itemize}
\item[(a)] $v(x)=1$ for every $x\in V_1$;
\item[(b)] $v(x)\ge1$ for every vertex $x\ne F$ visited by the walk;
\item[(c)] if $N\notin V_1$, every $x\in V_1$ has a neighbour $z$ with $w_{xz}>0$ and $v(z)>1$, and every vertex
with $v>1$ is joined to $N$ by a path of the support along which $v>1$;
\item[(d)] if $N\in V_1$, then $w_{NF}=1$, $v\equiv1$ on the visited vertices, and $w_{xF}=0$ for every other
neighbour $x$ of $F$.
\end{itemize}
\end{proposition}
\begin{proof}
(a) By Lemma~\ref{lem:cross}, $w_{xF}=p_{xF}(w)=w_{xF}v(x)$.
(b) By \eqref{eq:kirchhoff}, at a visited vertex outside $V_1\cup\{F\}$ the value $v(x)$ is at least the weighted mean
of $v$ over its neighbours in the support, none of which is $F$. If the minimum of $v$ over the visited vertices other
than $F$ were attained outside $V_1$, it would be attained at all support neighbours, and, the support component of
$N$ being connected and containing a vertex of $V_1$, also on $V_1$, where $v=1$.
(c) For $x\in V_1$, $x\neq N$, \eqref{eq:kirchhoff} and (a) give $\sum_{z\ne F}w_{xz}(v(z)-1)=w_{xF}>0$, so some
neighbour $z$ in the support has $v(z)>1$. Now let $v(z)=t>1$ and let $C$ be the set of vertices joined to $z$ by a
path of the support along which $v\ge t$. Suppose $N\notin C$. By (a), $C$ does not meet $V_1$, so by
\eqref{eq:kirchhoff} the value of $v$ at a vertex of $C$ is the weighted mean of $v$ over its neighbours in the
support. A neighbour outside $C$ has $v<t$. Hence a vertex where $v$ is maximal on $C$ has all its support neighbours
in $C$, and $v$ is maximal at each of them. Repeating, $C$ is a union of connected components of the support. But $z$
is visited, so its component contains $N$, a contradiction.
(d) If $N\in V_1$ then $v(N)=1$, and $v(N)$ is the maximum of $v$, so $v\equiv1$ on the visited vertices by (b).
Then \eqref{eq:kirchhoff} at $N$ reads $w_{NF}=1$, and at another $x\in V_1$ it would read $w_{xF}=0$.
\end{proof}

\begin{proof}[Proof of Theorem~\ref{thm:shadow}]
Let $w$ be an equilibrium with all coordinates positive. Then $V_1=N(F)$, the support is $E$, and by
Proposition~\ref{prop:volt}(a) the set $\{v>1\}$ does not meet $N(F)$. By (c), each $x\in N(F)$ has a neighbour $z$
with $v(z)>1$, and $z$ is joined to $N$ by a path inside $\{v>1\}\subset V\setminus(N(F)\cup\{F\})$; thus $z\in I_0$.

For the second statement, $W_e(n)-\sum_{j<n}p_e(W(j))$ is a martingale with increments bounded by $1$, so
$n^{-1}\bigl(W_e(n)-\sum_{j<n}p_e(W(j))\bigr)\to0$ almost surely. On the event that $W(n)/n\to\chi$ with all
$\chi_e>0$, homogeneity and the continuity of $p$ on $(0,\infty)^E$ give $p_e(W(j))\to p_e(\chi)$, hence
$p(\chi)=\chi$, which has just been excluded.
\end{proof}

Lemma~\ref{lem:cross} also explains which zero edges can be stable. Let $w$ be an equilibrium and $e=\{x,z\}$ an
edge with $w_e=0$. If a small weight $\varepsilon$ is put on $e$ alone, then $p_e=\varepsilon\,(v(x)+v(z))+o(\varepsilon)$
when both endpoints are visited, and $p_e=\varepsilon\,v(x)+o(\varepsilon)$ when only $x$ is. By
Proposition~\ref{prop:volt}(b) the first rate is at least $2$ unless $z=F$, so such an edge is repelled from zero. An
edge to an unvisited vertex, or to $F$, is not repelled only when $v(x)=1$, and then the rate equals $1$ exactly:
linear analysis is silent, the decay is of order $1/t$, and whether the edge dies is decided at second order, as in
Theorem~\ref{thm:dead}.

\section{A corrected conjecture and a census of small graphs}\label{sec:census}

Let $(G,N,F)$ be reduced. If $N$ is adjacent to $F$ put $I_0=\emptyset$ and $A=\{N\}$; otherwise let $I_0$ be the
component of $N$ in $G-F-N(F)$ and let $A$ be the set of neighbours of $F$ that have a neighbour in $I_0$. Call
$V\setminus(I_0\cup\{F\})$ the \emph{shadow}. Let $\mathcal D(G,N,F)$ be the set of edges of the following two kinds:
\begin{itemize}
\item[(D1)] $\{x,F\}$ with $x\in N(F)\setminus A$;
\item[(D2)] for every shadow vertex $u$ and every bridge $e=\{u,y\}$ of $G-F$ whose far side $H$ (the component of $y$
in $G-F-e$) does not contain $N$: the edge $e$ and every edge with an endpoint in $H$.
\end{itemize}

\begin{conjecture}\label{conj:new}
For every reduced $(G,N,F)$, $W(n)/n$ converges almost surely to a deterministic limit $\chi$, and $\chi_e=0$ if and
only if $e\in\mathcal D(G,N,F)$.
\end{conjecture}

Theorem~\ref{thm:dead} proves $\chi_e=0$ for the edges of kind (D2) with $u\in N(F)$. When $N$ is adjacent to $F$,
the edges of kind (D1) are covered by \cite[Proposition 2.8]{KMS}, and Proposition~\ref{prop:volt}(d) is the
corresponding statement for equilibria. When $N$ is not adjacent to $F$, Theorem~\ref{thm:shadow} shows that
$\mathcal D\neq\emptyset$ is forced as soon as (D1) is non-empty.

We tested the mean-field form of Conjecture~\ref{conj:new} on every connected graph with $3$ to $7$ vertices and
every ordered pair $(N,F)$ up to automorphism, $22{,}828$ cases in all, of which $20{,}092$ are reduced; in the others the
never-crossed edges were removed first. For each case the flow $\dot w=p(w)-w$ was integrated from six initial
conditions (uniform, log-uniform over six $e$-folds, and concentrated on a random path from $N$ to $F$); coordinates
that decay were set to zero and the remaining ones were refined by Newton's method, with $p$ evaluated by one linear
solve per edge.
\begin{center}
\begin{tabular}{lrrrrr}
\toprule
 & cases & one limit & zero set $=\mathcal D$ & reduced & reduced, $\mathcal D\neq\emptyset$\\
\midrule
$N$ adjacent to $F$ & 12{,}004 & 12{,}003 & 12{,}003 & 10{,}286 & 10{,}064\\
distance at least $2$ & 10{,}824 & 10{,}824 & 10{,}824 & 9{,}806 & 4{,}118\\
\bottomrule
\end{tabular}
\end{center}
The six runs reached the same point in every case but one. In the remaining case, a graph with seven vertices in which
$N$ is adjacent to $F$, an edge of $\mathcal D$ has a weight that decays like $1/t$ along the flow, and in one run it had
not been set to zero by the end of the integration; with ten times the integration time the six runs agree and the
zero set is $\mathcal D$. At the limits the largest real part of an eigenvalue of
$Dp-I$ restricted to the support is $-0.015$, and the smallest positive coordinate is $0.0019$. Among the $9{,}806$
reduced cases at distance at least $2$, Theorem~\ref{thm:dead} applies to $2{,}504$ and Theorem~\ref{thm:shadow} to
$3{,}294$; $1{,}643$ of the $4{,}118$ cases with a dead edge are $2$-connected graphs, and $247$ are $3$-connected.

The process itself was simulated with $10^7$ ants on four graphs. On the graph of Example~\ref{ex:c4} the values of
$W_{ay}(n)/n$ at $n=10^3,10^4,\dots,10^7$ were $0.060$, $0.059$, $0.049$, $0.044$, $0.040$, and those of
$W_{yF}(n)/n$ were $0.020$, $0.016$, $0.010$, $0.007$, $0.005$, while the four edges of the cycle were within $0.02$
of their limits $2^{-1/2}$ and $1/2$ at $n=10^7$. On the two graphs described after Theorem~\ref{thm:shadow} the
weight of the edge at $F$ that dies in the flow went from $0.123$ to $0.067$ and from $0.098$ to $0.044$ over the same
range. The decay is slow. In Step~2 of the proof of Theorem~\ref{thm:dead} the logarithm of $X_n/A_n$ decreases in
conditional mean by an amount of order $p_e/n$ per ant, and $p_e$ is itself of order $X_n/n$; this points to a rate of
order $1/\log n$, which is compatible with these values.
\section{Theta graphs}\label{sec:theta}

Let $\Theta(L_1,\dots,L_k)$ consist of $k\ge2$ internally disjoint paths from $N$ to $F$, the $i$-th of which has
$L_i\ge2$ edges. For $k=2$ the almost sure limit of $W(n)/n$ was found in \cite[Theorem 1.5]{KMS}: the $j$-th edge
from $N$ of each path has limit $\alpha^j$, with a ratio $\alpha$ depending on the path. The same form gives an
equilibrium for every $k$.

\begin{proposition}\label{prop:theta}
There is a unique $(\alpha_1,\dots,\alpha_k)\in(0,1)^k$ such that
\begin{equation}\label{eq:theta}
\sum_{i=1}^k\alpha_i^{L_i}=1\qquad\text{and}\qquad\frac{1-\alpha_i}{1-\alpha_i^{L_i}}\ \text{does not depend on }i,
\end{equation}
and the vector that gives weight $\alpha_i^{\,j}$ to the $j$-th edge from $N$ of the $i$-th path, $1\le j\le L_i$, is
an equilibrium on $\Theta(L_1,\dots,L_k)$. If $L_1=\dots=L_k=L$ then $\alpha_i=k^{-1/L}$ for all $i$.
\end{proposition}
\begin{proof}
Let the $j$-th edge of path $i$ have weight $\alpha_i^{\,j}$, with $\alpha_i\in(0,1)$ to be chosen. This edge is
crossed exactly when the walk reaches the vertex at depth $j$ of path $i$. An excursion from $N$ enters path $i$ with
probability proportional to $\alpha_i$, and then reaches depth $j$ before returning to $N$ with probability
$\bigl(\sum_{m=0}^{j-1}\alpha_i^{-m}\bigr)^{-1}$, by the gambler's ruin formula for the path with resistances
$\alpha_i^{-1},\dots,\alpha_i^{-j}$. So the excursions that reach depth $j$ of path $i$ occur at rate
\[
a_{i,j}=\frac{\alpha_i}{\sum_{m=0}^{j-1}\alpha_i^{-m}}=\frac{1-\alpha_i}{\alpha_i^{-j}-1},
\]
and those that reach $F$ along path $i$ at rate $c_i=a_{i,L_i}$. The walk ends at the first excursion that reaches
$F$, so with $C=\sum_ic_i$,
\[
p_{i,j}=\frac{a_{i,j}}{a_{i,j}+C-c_i}.
\]
Since $a_{i,j}(1-\alpha_i^{\,j})=(1-\alpha_i)\alpha_i^{\,j}$, the equation $p_{i,j}=\alpha_i^{\,j}$ is equivalent to
$1-\alpha_i=C-c_i$, for every $j$. As $c_i=(1-\alpha_i)\alpha_i^{L_i}/(1-\alpha_i^{L_i})$, this reads
$(1-\alpha_i)/(1-\alpha_i^{L_i})=C$, and then $c_i=C\alpha_i^{L_i}$, so that $C=\sum_ic_i$ becomes
$\sum_i\alpha_i^{L_i}=1$. This is \eqref{eq:theta}.

The function $\alpha\mapsto(1-\alpha)/(1-\alpha^{L})=(1+\alpha+\dots+\alpha^{L-1})^{-1}$ decreases from $1$ to $1/L$
on $(0,1)$. For $C\in(\max_i1/L_i,1)$ let $\alpha_i(C)$ be the solution of $(1-\alpha)/(1-\alpha^{L_i})=C$. Then
$\sum_i\alpha_i(C)^{L_i}$ is continuous and decreasing in $C$, tends to $0$ as $C\to1$, and exceeds $1$ as
$C\to\max_i1/L_i$, because one term tends to $1$ and the others stay positive. So \eqref{eq:theta} has exactly one
solution. In the case of equal lengths the second condition gives equal $\alpha_i$ and the first gives
$\alpha^L=1/k$.
\end{proof}

For $k=2$, \eqref{eq:theta} is the system of \cite[Theorem 1.5]{KMS}. In the flow, thirty random initial conditions
converged to the equilibrium of Proposition~\ref{prop:theta} on each of $39$ theta graphs with $k\le8$ and
$L_i\le8$.

\section*{Acknowledgements and funding}
This work was done by multiple AI agents, Claude (Anthropic) and Codex (OpenAI), connected to Math
Brain~\cite{MathBrainRepo} and working together inside the Brain framework~\cite{BrainFramework}, a shared, persistent
memory in which the agents record claims, proofs, failed attempts and methods, retrieve one another's results,
coordinate who works on what, and audit each other's work. For more on the details behind the Brain and Math Brain, see
the Math Brain paper~\cite{MathBrainPaper}, the Math Brain repository~\cite{MathBrainRepo} and the Brain framework
repository~\cite{BrainFramework}.

The research presented in this paper was supported by the European Research Council (ERC) under the European Union's
Horizon Europe research and innovation programme (grant agreement No.~101041711), by the Simons Foundation as part of
the Collaboration on the Mathematical and Scientific Foundations of Deep Learning, by Heights Labs, by Convex Nexus
Capital, by the Israel Science Foundation (grant number 2258/19), and by the Israel Science Foundation (ISF Grant
4101/25). Computational resources were provided in part by Google Cloud through its education credits program (credit
530550103).

\section*{Declaration of competing interest}
The author holds equity in Heights Labs and in Convex Nexus Capital and has a management role at Convex Nexus Capital.
These entities are included in the funding declaration above.

\section*{Declaration of generative AI and AI-assisted technologies}
The paper comes out of Math Brain, an ongoing project on the use of AI systems for research mathematics, and its
results were developed with the Brain~\cite{BrainFramework} and Math Brain~\cite{MathBrainRepo}. In November 2026 a
Claude Fable agent (Anthropic), connected to the Brain and working from the project's list of open problems in
probability, found the statements of this paper and their proofs, wrote the programs, ran the census and the
simulations, and prepared the manuscript. A second Claude agent independently checked part of the computations.
These checks are not peer review. The author directed the work and is responsible for the content of the paper and for
decisions concerning its dissemination.

\section*{Code and data availability}
Programs, census files and outputs:\\
\url{https://github.com/yspennstate/trace-reinforced-ants-dead-branches}.

The paper is part of the Math Brain project, whose public repository is
\url{https://github.com/yspennstate/math-brain-public}~\cite{MathBrainRepo}.

\begin{thebibliography}{99}
\bibitem{Ben} M.~Bena\"im, Dynamics of stochastic approximation algorithms, in \emph{S\'eminaire de Probabilit\'es
XXXIII}, Lecture Notes in Math. 1709, Springer, 1999, 1--68.
\bibitem{Dav} B.~Davis, Reinforced random walk, \emph{Probab. Theory Related Fields} \textbf{84} (1990), 203--229.
\bibitem{Den} J.-L.~Deneubourg, S.~Aron, S.~Goss and J.~M.~Pasteels, The self-organizing exploratory pattern of the
Argentine ant, \emph{J. Insect Behav.} \textbf{3} (1990), 159--168.
\bibitem{DS} M.~Dorigo and T.~St\"utzle, \emph{Ant Colony Optimization}, MIT Press, 2004.
\bibitem{Pem90} R.~Pemantle, Nonconvergence to unstable points in urn models and stochastic approximations,
\emph{Ann. Probab.} \textbf{18} (1990), 698--712.
\bibitem{Pem07} R.~Pemantle, A survey of random processes with reinforcement, \emph{Probab. Surv.} \textbf{4}
(2007), 1--79.
\bibitem{KMS0} D.~Kious, C.~Mailler and B.~Schapira, Finding geodesics on graphs using reinforcement learning,
arXiv:2010.04820.
\bibitem{KMS} D.~Kious, C.~Mailler and B.~Schapira, The trace-reinforced ants process does not find shortest paths,
\emph{J. \'Ec. polytech. Math.} \textbf{9} (2022), 505--536.
\bibitem{MV} C.~Mailler and Z.~Varin, The two-nest ants process on triangle-series-parallel graphs,
arXiv:2601.22855.
\bibitem{MathBrainRepo} Y.~Shmalo, \emph{Math Brain: public snapshot, code and tests}, 2026,
\url{https://github.com/yspennstate/math-brain-public}.
\bibitem{MathBrainPaper} Y.~Shmalo, \emph{Math Brain: shared, auditable memory for multi-agent mathematical research},
October 2026, \url{https://github.com/yspennstate/math-brain-public/tree/main/papers/math-brain-2026-10-24}.
\bibitem{BrainFramework} Y.~Shmalo, \emph{Brain Framework}, public software repository, 2026,
\url{https://github.com/yspennstate/brain-framework-public}.
\end{thebibliography}
\end{document}
